\section{Introduction}

Reinforcement learning has become an important component of large language model (LLM) training pipelines, particularly for alignment, preference optimization, and reasoning tasks. Most contemporary approaches rely on policy gradient methods augmented with advantage estimation, typically implemented via learned value functions or generalized advantage estimation. These techniques originate from classical reinforcement learning settings in which states are well-defined, approximately Markov, and admit stable value estimates.

Applying this framework to language generation introduces a fundamental mismatch. In autoregressive LLMs, the \``state\'' is a partial sequence of tokens. Such prefixes are not Markov, do not summarize progress toward task completion, and can vary widely in semantic meaning despite superficial similarity. Consequently, estimating a value function over prefixes is both statistically challenging and conceptually ill-defined. Empirically, value heads in language RL are often brittle, sensitive to distributional shift, and prone to introducing bias that outweighs their variance-reduction benefits.

Motivated by these issues, recent work has increasingly moved away from learned critics in LLM reinforcement learning. Critic-free, outcome-based methods such as RLOO and GRPO replace value estimation with batch-level baselines, achieving stable training on verifiable reasoning tasks such as math and code. However, these approaches still implicitly assume uniform credit assignment across tokens: every action in a trajectory receives the same reward signal, modulated only by a scalar baseline.

This raises a natural question: \emph{if we no longer estimate values over prefixes, how should trajectory-level reward be distributed across tokens?} Uniform credit assignment is simple, but it ignores the fact that different tokens play very different roles in generation. Some tokens correspond to high-uncertainty decisions, strategy commitments, or branching points in reasoning, while others are low-entropy continuations or stylistic filler. Treating all tokens equally risks diluting the learning signal and obscuring the mechanisms by which policies improve.

In this work, we study token-level credit redistribution as a mechanism for on-policy LLM reinforcement learning without value functions. Rather than estimating advantages or prefix values, we propose to weight token-level policy gradients using intrinsic signals derived directly from the model's predictive distribution. These include token surprisal and changes in predictive entropy, which naturally reflect decision difficulty and commitment during generation. Importantly, these signals are inexpensive to compute, require no additional networks or tree structures, and make no assumptions about the existence of a meaningful value function over prefixes.

To make this question concrete, the repository now implements a reproducible experimental stack around Qwen models on GSM8K and MBPP. The code compares uniform token updates against RLOO/GRPO baselines with surprisal-based and entropy-reduction-based intrinsic weighting, all under matched prompt formatting and optimization settings. This turns the paper's experimental section into a directly runnable benchmark rather than a placeholder description.

Overall, this paper reframes advantage estimation as one instantiation of a broader credit assignment problem in language reinforcement learning. In settings where state semantics are weak and rewards are trajectory-level, redistributing credit using intrinsic distributional signals is a plausible alternative to value-based methods; the implemented experiments are designed to test that claim directly.
